\documentclass[12pt,a4paper]{article}

% --- Packages ---
\usepackage[utf8]{inputenc}
\usepackage[T1]{fontenc}
\usepackage[french]{babel}
\usepackage{geometry}
\geometry{margin=2.5cm}
\usepackage{graphicx}
\usepackage{hyperref}
\usepackage{enumitem}
\usepackage{booktabs}
\usepackage{xcolor}
\usepackage{titlesec}
\usepackage{fancyhdr}
\usepackage{float}

% --- Configuration de la mise en page ---
\pagestyle{fancy}
\fancyhf{}
\lhead{Rapport d'Avancement - SpeechCoach}
\rhead{Master SIE - 2026}
\cfoot{\thepage}

\titleformat{\section}{\large\bfseries\color{blue!70!black}}{}{0em}{}[\titlerule]

% --- Informations du document ---
\title{
    \vspace{-2cm}
    \large Master Systèmes Intelligents pour l'Éducation (SIE)\\
    \vspace{0.5cm}
    \huge \textbf{Rapport d'État d'Avancement}\\
    \large \textit{Projet : SpeechCoach}
}
\author{\textbf{Étudiant :} PFE Student (Hisham Moussaid) \\ \textbf{Encadrant :} Pr. Abdelaaziz Hessane}
\date{Période du 26 Janvier 2026 au 20 Février 2026}

\begin{document}

\maketitle

\section{Synthèse du Jalon}
\begin{itemize}
    \item \textbf{Jalon visé :} Fin Sprint 5 (Initialement prévu mi-avril) - Pose / Gestes (niveau MVP) et socle multimodal.
    \item \textbf{État global :} \textbf{En avance}. Les jalons 0 à 5 ont été réalisés avec succès sur une période d'un mois.
    \item \textbf{Objectif principal de la période :} Mettre en place la fondation technique (ingestion vidéo) et consolider les briques d'analyse audio (ASR, rythme, pauses) ainsi que les premières briques de vision (visage, eye contact, activité des mains).
\end{itemize}

\section{Travaux Réalisés \& Résultats Techniques}
\subsection{Implémentation des Modules}
\begin{itemize}
    \item \textbf{Code développé :} Développement du script d'orchestration \texttt{process\_video.py} gérant l'ingestion multimodale, ainsi que des modules \texttt{audio/asr.py}, \texttt{audio/analytics.py} et \texttt{vision/analysis.py}.
    \item \textbf{Modèles intégrés :} Intégration de \textbf{faster-whisper} (int8) pour une transcription rapide et précise de la parole, et de \textbf{MediaPipe} (Face Mesh, Hands, Pose) pour le tracking facial, l'axe du regard et l'intensité gestuelle.
    \item \textbf{Extraction de métriques :}
    \begin{itemize}
        \item \textit{Audio} : Débit de parole (WPM), détection de pauses ($>0.5\text{s}$) via l'énergie RMS (Librosa), comptage de fillers.
        \item \textit{Vision} : Ratio de présence visage, vérification de luminosité/flou, estimation de l'Eye Contact (Yaw/Pitch estimé via solvePnP), ratio de visibilité des mains et intensité des mouvements corporels.
    \end{itemize}
\end{itemize}

\subsection{Performance et Validation (DoD)}
Afin de valider le système (DoD respectées pour Sprint 0-5), plusieurs tests batch ont été effectués :
\begin{itemize}
    \item \textbf{Stabilité :} Le système sort de manière constante un objet JSON stéréotypé contenant les données brutes (\texttt{report.json}), validant le contrat de l'architecture.
    \item \textbf{Temps de traitement :} Sur une architecture CPU, on obtient actuellement des performances d'inférence allant de \textbf{2.48x à 2.77x le RTF} (Real Time Factor), soit ~10 minutes de traitement pour ~4 minutes de vidéo.
    \item \textbf{Livrables produits :} Rapports d'analyse combinés au format Markdown (générant un "feedback" initial), graphiques `audio\_energy.png` (timeline énergie et silences) et `vision\_timeline.png` (timeline vision et tracking), plus les transcripts bruts.
\end{itemize}

\subsubsection*{Échantillons de validation récents :}
\begin{itemize}
    \item \textbf{Test 2 (Rapide) :} 194 WPM, Eye Contact 54\%, Intensité gestuelle 2.8/10 (corps figé avec présence de mains 9\% du temps). Traitement RTF 2.48x.
    \item \textbf{Test 5 (Posé) :} 143 WPM, Eye Contact 86\%, Intensité gestuelle 3.3/10 (présence de mains 94\% du temps). Traitement RTF 2.77x.
\end{itemize}

\section{Stack Technique et Frameworks}
\begin{table}[H]
    \centering
    \begin{tabular}{@{}ll@{}}
        \toprule
        \textbf{Composante} & \textbf{Outils / Librairies} \\
        \midrule
        Traitement Vidéo & FFmpeg, OpenCV \\
        Transcription \& ASR & Faster-Whisper (modèle "medium" int8) \\
        Analyse Audio & Librosa, Numpy \\
        Analyse Vision & MediaPipe (Face Mesh, Pose, Hands) \\
        Visualisation & Matplotlib \\
        Output \& Serialization & Types natifs Python, JSON, Markdown \\
        \bottomrule
    \end{tabular}
    \caption{Environnement de développement technique validé pour les Sprints 0-5.}
\end{table}

\section{Obstacles et Défis Rencontrés}
\begin{itemize}
    \item \textbf{Défis Techniques \& Dépendances :}
    \begin{itemize}
        \item \textbf{Conflits de bibliothèques} : Le passage à NumPy 2.x cassait la compatibilité binaire de MediaPipe et générait des erreurs d'attributs. Il a fallu créer un script de résolution (\texttt{fix\_dependencies.py}) pour forcer le downgrade (NumPy < 2.0.0, Protobuf 3.20.3) et stabiliser l'environnement.
        \item Le comportement de Whisper "aseptise" certaines hésitations ("euh", "hmm"), faussant le décompte des tics de langage.
        \item L'angle par défaut des webcams de PC portable force un regard subjectif vers le bas, faussant la métrique de l'Eye Contact pur.
        \item Les seuils de silence absolus sont faussés par le bruit de fond, nécessitant un seuil relatif adaptatif (\texttt{top\_db=30}).
    \end{itemize}
    \item \textbf{Contraintes Matérielles :} L'inférence en parallèle de Whisper avec MediaPipe sur un environnement purement CPU occasionne un délai un peu supérieur au temps réel (RTF $\approx 2.5$).
    \item \textbf{Solutions :} Asymétrie appliquée aux marges d'angle Pitch ($-25^\circ$). Adaptation algorithmique pour le filtrage audio. Verrouillage strict des versions de l'environnement Python.
\end{itemize}

\section{Planning et Prochaines Étapes}
\begin{itemize}
    \item \textbf{Retard éventuel :} Aucun retard, le projet avance de manière exceptionnellement rapide par rapport à la macro-programmation (Sprints prévus jusqu'en avril terminés mi-février).
    \item \textbf{Objectifs du prochain sprint (Sprint 6) :} \textbf{Agrégation et scoring pédagogique}. Le but est d'utiliser toutes ces métriques pour calculer des notes sur des axes précis (ex. Rythme, Regard, Gestuelle) et de générer un résumé déterministe (sans LLM) listant "Forces" et "Faiblesse".
    \item \textbf{Priorités (Plan d'accélération) :} Focus critique sur une rubrique d'évaluation pédagogique correcte et justifiée qui servira de base solide avant d'aborder les générations de conseils et la phase RAG (Sprints 7 \& 8).
\end{itemize}

\section{Références et Veille}
\begin{itemize}
    \item Documentation officielle de \textbf{Google MediaPipe Solutions} (Face Mesh, Pose, Hands).
    \item \textit{Robust Speech Recognition via Large-Scale Weak Supervision} (Radford et al., 2022) - Modèle \textbf{OpenAI Whisper}.
    \item \textbf{CTranslate2} (Système d'inférence) via l'implémentation \texttt{faster-whisper}.
    \item Bibliothèque d'analyse audio \textbf{Librosa} (McFee et al., 2015) pour l'extraction de l'énergie (RMS).
    \item \textbf{OpenCV} (Open Source Computer Vision Library) pour le calcul matriciel de head pose (\texttt{solvePnP}) et la détection bas niveau.
    \item Modèle tridimensionnel (Mehrabian's Communication Model) justifiant la pondération de l'impact du langage non-verbal face au verbal dans la qualité oratoire.
\end{itemize}

\end{document}
